\documentclass{article}
\usepackage{amsmath}
\usepackage{amssymb}
\usepackage{amsthm}
\usepackage{mathtools}
\usepackage{geometry}
\usepackage{hyperref}

\newcommand{\rk}{\operatorname{rank}}
\newcommand{\GL}{\operatorname{GL}}
\newcommand{\R}{\mathbb{R}}
\newcommand{\Mat}{\mathrm{Mat}}

\theoremstyle{definition}
\newtheorem{theorem}{Theorem}[section]
\newtheorem{proposition}[theorem]{Proposition}
\newtheorem{lemma}[theorem]{Lemma}
\newtheorem{corollary}[theorem]{Corollary}
\newtheorem{definition}[theorem]{Definition}
\newtheorem{remark}[theorem]{Remark}
\newtheorem{example}[theorem]{Example}

\title{Examples and Counterchecks for the Second Conjecture}
\date{}

\begin{document}
\maketitle

\section{Purpose}

This file records the low-dimensional checks that validate the setup of the
second-conjecture block.

The checks to keep visible are:

\begin{itemize}
\item the path $a_L\cdots a_1$ survives;
\item the relations kill exactly the length-$L$ paths through the closing arrow;
\item the Kupisch series is $(L+1,L,\dots,L)$;
\item the admissible intervals split into a non-wrapping type-$A$ family and a
  wrapped family.
\end{itemize}

\section{The case $L=2$}

\begin{example}[The algebra $A_2^{\mathrm{wt}}$]
The quiver is
\[
0 \xrightarrow{a_1} 1 \xrightarrow{a_2} 2 \xrightarrow{a_0} 0.
\]
The weight-space relations are
\[
p_1=a_0a_2=0,
\qquad
p_2=a_1a_0=0.
\]
The surviving length-$2$ path is
\[
p_0=a_2a_1,
\]
which is the quiver version of the end-to-end map.
\end{example}

\begin{proposition}[Projectives for $L=2$]
The projective lengths are
\[
(c_0,c_1,c_2)=(3,2,2).
\]
\end{proposition}

\begin{proof}
Starting at $0$, the path $a_2a_1$ survives, so the maximal nonzero path length
is $2$.
Starting at $1$ or $2$, the unique length-$2$ path goes through the closing
arrow and is one of the generators, so the maximal nonzero path length is $1$.
\end{proof}

\begin{remark}[Admissible intervals for $L=2$]
The non-wrapping family is
\[
M(0,0),\ M(0,1),\ M(0,2),\ M(1,1),\ M(1,2),\ M(2,2).
\]
The wrapped family is the single module
\[
M(2,0).
\]
The forbidden wrapped interval is
\[
M(1,0),
\]
because it would have cyclic length $2$ through the closing arrow.
\end{remark}

\section{The case $L=3$}

\begin{example}[The algebra $A_3^{\mathrm{wt}}$]
The quiver is
\[
0 \xrightarrow{a_1} 1 \xrightarrow{a_2} 2 \xrightarrow{a_3} 3 \xrightarrow{a_0} 0.
\]
The relations are
\[
p_1=a_0a_3a_2=0,
\qquad
p_2=a_1a_0a_3=0,
\qquad
p_3=a_2a_1a_0=0.
\]
The surviving length-$3$ path is
\[
p_0=a_3a_2a_1.
\]
\end{example}

\begin{remark}[Comparison with the gradient equations]
Under the identification
\[
a_0\rightsquigarrow W_0,
\qquad
a_\ell\rightsquigarrow W_\ell,
\]
the above three relations become exactly
\[
W_0W_3W_2=0,
\qquad
W_1W_0W_3=0,
\qquad
W_2W_1W_0=0,
\]
which is the $L=3$ example from
\texttt{gradient\_equations\_in\_weight\_space.tex}.

The open-chain product
\[
W_3W_2W_1
=
\mu(W)
\]
is not killed.
\end{remark}

\begin{proposition}[Projectives for $L=3$]
The Kupisch series is
\[
(c_0,c_1,c_2,c_3)=(4,3,3,3).
\]
\end{proposition}

\begin{remark}[Admissible intervals for $L=3$]
The non-wrapping family is
\[
M(i,j)
\qquad
(0\le i\le j\le 3).
\]
The wrapped family is
\[
M(2,0),\qquad M(3,0),\qquad M(3,1).
\]
The forbidden wrapped intervals are
\[
M(1,0),\qquad M(2,1),\qquad M(3,2).
\]
\end{remark}

\section{Edge checks to keep in mind}

\begin{remark}[Allowed versus forbidden length-$L$ intervals]
The open-chain interval
\[
M(0,L)
\]
is allowed because it comes from the surviving path
\[
a_L\cdots a_1.
\]

By contrast, for each
\[
1\le i\le L,
\]
the wrapped interval
\[
M(i,i-1)
\]
is forbidden because it would require the vanishing path $p_i$.
\end{remark}

\begin{remark}[Why the module language matches the cyclic locus]
After fixing the closing arrow to
\[
a_0=W_0=\Sigma_{XY}P_Q,
\]
the defining quiver relations become exactly the equations
\[
W_{<\ell}W_0W_{>\ell}=0.
\]
Since the path
\[
a_L\cdots a_1
\]
survives, this reformulation does not accidentally impose
\[
\mu(W)=0.
\]
\end{remark}

\section{Deliverables}

This file is complete once it records:

\begin{itemize}
\item the explicit `L=2` and `L=3` relations;
\item the explicit projective lengths in those cases;
\item the explicit admissible interval lists;
\item the allowed open-chain length-$L$ interval and the forbidden wrapped
  length-$L$ intervals.
\end{itemize}

\end{document}
